\documentclass[a4paper, 12pt, openany]{book} %chose the paper size and font size. Openany ensures that all all chapters and similar may begin at any page, not only odd pages. For the introductory pages and appendices we want openany, but for chapter pages in the main content we want chapters to begin only on odd pages (right hand side). The book class ensures that the margins are automatically adjusted such that left hand pages are slightly moved to the left and vice versa at the right, which makes the thesis very readable and good looking when printed in bound book format.
\usepackage[utf8]{inputenc} %to manage special characters
\usepackage[T1]{fontenc} %to manage special characters
\usepackage[Bjarne]{fncychap} %fancy chapter style (many more available, like Sonny or Lenny etc.)
\usepackage{fancyhdr} %to customize the headers
\usepackage[lmargin=1.5in, rmargin=1in, tmargin=1in, bmargin=1in]{geometry} %sets the margins for the pages
\setcounter{tocdepth}{2} %table of contents number depth for subsections (2 = x.x.x)
\setcounter{secnumdepth}{4} %numbering depth for headers for subsections in the text(4 = x.x.x.x)
\usepackage{url} %to include urls
\usepackage{listings} %include this if you want to include code in the thesis
\usepackage{amsmath,amssymb} %mathematical package
\usepackage{siunitx} %includes SI-units
\usepackage[bf]{caption} %makes float captions bold
\usepackage{array, booktabs} %to make better tables
\usepackage{graphicx} %to include graphics
\usepackage{float} %to include floats
\usepackage[export]{adjustbox} %to adjust floats
\usepackage{subfig} %to include subfigures
\usepackage{chngcntr} %will make it possible to change the counter for tables, figures etc. such as below
\counterwithin{figure}{section} %change counter for figures within sections (also possible to choose for each chapter
\counterwithin{table}{section} %change counter for tables within sections
\usepackage{color, xcolor} %edit e.g. text colors
% Importing document settings from our file "packages.sty"
\usepackage{packages}
\usepackage{tikz}
\usetikzlibrary{positioning}
\usepackage{subfig}
\usepackage{comment}


\usepackage[backend = biber,
            style = numeric,
            date = long,     % Long: 24th Mar. 1997 | Short: 24/03/1997
            sorting = none,
            maxcitenames = 3,   % max names to include before et. al.
            ]{biblatex} %customize the look of your citations and bibliography
%\addbibresource{references.bib} %declare the bibliography resource
\usepackage{comment} %to be able to comment out sections in the .tex files
\usepackage{afterpage} %to customize page commands such as below
\newcommand\myemptypage{
    \null
    \thispagestyle{empty}
    \addtocounter{page}{-1}
    \newpage
    } %sets new page command to insert an empty page without adding to the page counter or having a page number




\begin{document}
%%%%%%%%%%%%%%%%%%%%%%%%%%%%%%%%%%%%%%%%%%%%%%%%%%%%%%%%
%\begin{comment}
% The title page:
% For NTNU students this page will be generated automatically when submitting your paper, and should not be included in the final file from Latex. Delete or comment out the title page setup. The final report should then start with the first page being the abstract. I have included a title page here so it is possible to see how it may look like, and for those who does not get an automatically generated title page. Of course you will need to change the names and titles etc. to your case.

%the title page should be an odd page (right hand side)

\begin{titlepage}
\newgeometry{left=1.6in, right=2in}
\vspace*{1.5cm}

\noindent  \textcolor{gray}{\large Martin Emil Pettersen} \\
\vspace{1cm}

\noindent \textbf{\Large The title of your master's thesis should be written here} \\
\vspace{0.5cm}

\noindent {\large Any undertitle is written here} \\



\vspace{7cm}
\noindent Master's thesis in Mathematical Sciences \\
Supervisor: Robert Brian O'Hara \\
Co-supervisor: Ron Ronald Togunov \\
June 2025 \\

\vspace{0.2cm}
\noindent Norwegian University of Science and Technology \\
Faculty of Information Technology and Electrical Engineering \\
Department of Mathematical Sciences \\

\begin{figure}[h]
    \includegraphics[width=0.28\textwidth]{figures/ntnu_basic.png}
\end{figure}
\end{titlepage}
\restoregeometry
\myemptypage %empty page such that the abstract starts at the first right hand side after the title page
%\end{comment}
%%%%%%%%%%%%%%%%%%%%%%%%%%%%%%%%%%%%%%%%%%%%%%%%%%%%%%%%


% The pre-chapters
\chapter*{Abstract} %pre-chapters should not be numbered, hence the "*"
\pagenumbering{roman} %introductory pages should be roman
\setcounter{page}{1}
\addcontentsline{toc}{chapter}{\protect\numberline{}Abstract} %add the chapter to the table of contents, this is not automatically added when creating unnumbered chapters (*). Add it in a chapter style, and keep all chapters on the same numberline indent regardless of number or not on the chapter


An animal's utilization distribution describes the probability of finding the animal at a given point in space. It is a tool that can be used to gain information about the habitat preferences of an animal. There exist many ways to model utilization distributions. The Langevin movement model is one model proposed for animal movement that does this. This model is based on the Langevin diffusion process. This process has an explicit stationary distribution that is the utilization distribution; thus, it can be used to model the long-term spatial distribution of an animal by looking at its displacement over time. To model data, the Langevin process is discretized using the Euler-Maruyama method. There is an error in this approximation and, for long time intervals between observations, this error leads to a bias in the utilization distribution parameter estimates. In this thesis, I study three methods for obtaining improved estimates for the Langevin movement model and find that using Monte Carlo integration over intermediate states between observations can be used to remove the bias. Using this method, the Langevin movement model can be applied to data where there are large intervals in time between observations. In addition, I show that when the sampling intervals are longer, there is a reduced variance in the utilization distribution parameter estimates. This means that the method can be used to improve animal tracking experiment designs by sampling at lower frequencies. 






\newpage
\null
\thispagestyle{empty}
\newpage

\chapter*{Sammendrag}













\newpage
\null
\thispagestyle{empty}
\newpage %insert the chapter text from the files

\chapter*{Preface}
\addcontentsline{toc}{chapter}{\protect\numberline{}Preface} 

Write the preface of your thesis here. \\

\noindent You may include acknowledgments and thanks as part of your preface on this page, or you may add it as a new chapter after the preface.









\section{Relevancy to UN development goals}
Studying animal habitats is an important part of species conservation. As we move into a future where the environment is increasingly affected by human pollution, studying animal behavior becomes more important to combat loss in biodiversity. The UN sustainability goals include goals to "combat climate change and its impacts", "Conserve and sustainably use the oceans, seas and marine resources for sustainable development" and "Protect, restore and promote sustainable use of terrestrial ecosystems, sustainably manage forests, combat desertification, and halt and reverse land degradation and halt biodiversity loss"\parencite{UN_goals}. This thesis is relevant to these goals, since it explores how to improve utilization distribution estimates. Improved utilization distribution estimates can be used to make informed decisions to ensure the future of endangered ecosystems.









\tableofcontents
\addcontentsline{toc}{chapter}{\protect\numberline{}Contents}

%add to table of contents list of figures and tables, and insert list of figures and tables
\addcontentsline{toc}{chapter}{\protect\numberline{}\listfigurename}
\listoffigures
\addcontentsline{toc}{chapter}{\protect\numberline{}\listtablename}
\listoftables


\chapter*{Abbreviations}
\addcontentsline{toc}{chapter}{\protect\numberline{}Abbreviations}
% Put in your abbreviations here

List of all abbreviations in alphabetic order:

\begin{itemize}
    \item \textbf{UD} Utilization Distribution
    \item \textbf{RSF} Resource Selection Function
    \item \textbf{Mamsl} meter above mean sea level
    \item \textbf{NTNU} Norwegian University of Science and Technology
    \item \textbf{PCA} Principal Component Analysis
    
\end{itemize}
\newpage
\myemptypage
%add an empty non-counted page by the command below in order to get the first chapter on the left hand side, if needed (check your page number so that the first chapter is on an odd page)


%%%%%%%%%%%%%%%%%%%%%%%%%%%%%%%%%%%%%%%%%%%%%%%%%%%%%%%%
%Customize the layout of the main content of your thesis

\pagestyle{fancy} %set customized page style for header
\fancyhf{} %clear header and footer fields
\renewcommand{\headrulewidth}{0pt} %set to no rule
\fancyhead[LE, RO]{\thepage} %set the page number at left for even, right for odd pages
\fancyhead[RE, LO]{\leftmark} %set the chapter name at right for even, left for odd pages
%is is possible to design the header with the chapter as you wish, e.q. only the chapter or only the name, all lowercase instead etc.
%you could also design the footer if you wish, for example:
%\fancyfoot[LE, RO]{\thepage}
\setlength{\headheight}{14.49998pt} %set the header height


%%%%%%%%%%%%%%%%%%%%%%%%%%%%%%%%%%%%%%%%%%%%%%%%%%%%%%%%
%main content 

\pagenumbering{arabic}
\chapter{Introduction}

A common problem in ecology is that of finding the home range of an animal. The home range was described by \parencite{burt1943territoriality} as "\dots that area traversed by the individual in its normal activities of food gathering, mating, and caring for young." This concept is formalized in statistics as the utilization distribution(UD), which \parencite{anderson1982home} defines as "the probability density function that gives the probability of finding an animal at a particular location". There exist many methods for estimating the UD of an animal, such as estimating it using occurrence data. In movement ecology, the UD is estimated using the movement of an animal in the form of track data. Many methods for modeling animal movement have been proposed; a model that is commonly used is the step selection analysis. In this thesis I focus on the Langevin movement model proposed in \parencite{michelot_langevin_2019}. This model uses the Langevin process to model 


A problem that \parencite{michelot_langevin_2019} observes is that for low-frequency sampling of animal movement, there is a bias in the parameters estimated by the Langevin movement model. This occurs because the discretization of the Langevin process used to model data is an approximate likelihood which does not take into account the movement that might have occurred between observations. 


This thesis explores three methods for dealing with the bias observed in \parencite{michelot_langevin_2019}. The first of these uses the extended Kalman filter(EKF), iterated to find predictions for the states of the process between observations, and a resulting likelihood approximation based on these predictions. The second method adds a grid of hidden states between observations. The transition densities between these hidden states are better approximations for the Langevin process, since they have shorter time intervals between them. An approximation for the Langevin likelihood is found by expressing the joint density of the observations and hidden states, and using the Monte-Carlo method to integrate out the hidden states. To improve the speed of this method, importance sampling is used. The third method uses the same concept as the second method, but rather than simulating proposals each time the likelihood is evaluated, an approximate proposal is simulated and used repeatedly instead. 


Of the methods described, the method that uses importance sampling to integrate over hidden states manages to reduce the bias seen in \parencite{michelot_langevin_2019}. This means that the Langevin movement model can be applied to cases where we have observations of animal movement sampled at a low frequency. In addition, when the observations have longer time intervals between them, there is a reduction in the variance of the estimated parameters of the UD. This means that an animal's UD can be better estimated if its position is gathered at lower frequencies.











\cleardoublepage
%the cleardoublepage command ensures that the next text page is on the right-hand side (odd page) and produces a blank page if necessary to achieve that, as all chapters should begin on the right hand side


% Defining front matter settings
\frontmatter

% Inserting table of contents
\tableofcontents

% Inserting list of figures & list of tables
\listoffigures

\listoftables

% Defining main matter settings (Norsk: innstillinger for hoveddelen av teksten)
\mainmatter

\chapter{Theory}
\import{./Sections/}{03Theory}
\cleardoublepage

\chapter{Background}
\import{./Sections/}{04Background}
\cleardoublepage

\chapter{Methods}
\import{./Sections/}{05Method}
\cleardoublepage

\chapter{Simulation studies}
\import{./Sections/}{06Results}
\cleardoublepage

\chapter{Discussion}
\import{./Sections/}{07Discussion}
\cleardoublepage

\chapter{Conclusion}
\import{./Sections/}{08Conclusion}
\cleardoublepage

% Printing bibliography
\newpage
\printbibliography[heading = bibintoc, title = Bibliography]    % 'bibintoc' inserts our bibliography into the table of contents
\cleardoublepage

% Inserting appendix with separate settings
\addappendix
\import{./Appendices/}{github}
\import{./Appendices/}{Appendix A}
\import{./Appendices/}{Appendix B}
\import{./Appendices/}{Appendix C}

% End of document
\end{document}
